\chapter{Descriptive Statistics}\label{ch:g10:stats}

Statistics summarizes a large collection of numbers --- grades, salaries,
temperatures --- by a few well-chosen indicators. This chapter introduces
the two families of indicators: those locating the \emph{center} of the
data (\omterm{def:g10:stats:mean}{mean}, \omterm{def:g10:stats:median}{median}) and those measuring its \emph{spread} (\omterm{def:g10:stats:quartiles}{range},
\omterm{def:g10:stats:quartiles}{quartiles}). A deeper treatment, including the standard deviation, comes in
\cref{ch:g11:stat}.

\section{Data and frequencies}

\begin{definition}[Statistical series]\label{def:g10:stats:series}
A \emph{statistical series}\index{statistical series} is a list of values
$x_1, x_2, \dots, x_N$ observed on a population of size $N$. When a value
$x$ appears $n$ times, $n$ is its \emph{count} and the quotient
\[
f = \frac{n}{N}
\]
is its \emph{frequency}\index{frequency} (often expressed as a
percentage). The frequencies of all values add up to $1$.
\end{definition}

\begin{example}\label{ex:g10:stats:frequencies}
The grades of $20$ students on a $5$-point quiz:
\begin{center}
\renewcommand{\arraystretch}{1.2}
\begin{tabular}{c|cccccc}
grade $x$ & $0$ & $1$ & $2$ & $3$ & $4$ & $5$ \\
\hline
count $n$ & $1$ & $2$ & $4$ & $6$ & $5$ & $2$ \\
\omterm{def:g10:stats:series}{frequency} $f$ & $0.05$ & $0.10$ & $0.20$ & $0.30$ & $0.25$ & $0.10$ \\
\end{tabular}
\end{center}
Check: $1 + 2 + 4 + 6 + 5 + 2 = 20$ and the frequencies sum to $1$.
\end{example}

\begin{omfigure}
\begin{tikzpicture}
\begin{axis}[
  width=8.4cm, height=5.2cm,
  ybar, bar width=14pt,
  xmin=-0.7, xmax=5.7, ymin=0, ymax=7.4,
  xtick={0,1,2,3,4,5},
  ytick={1,2,3,4,5,6},
  xlabel={grade}, ylabel={count},
  tick label style={font=\small},
  label style={font=\small}]
  \addplot[fill=omDef!30, draw=omDef] coordinates
    {(0,1) (1,2) (2,4) (3,6) (4,5) (5,2)};
\end{axis}
\end{tikzpicture}

{\small The bar chart of the quiz grades: one bar per value, with height
equal to the count.}
\end{omfigure}

\section{Measures of center}

\begin{definition}[Mean]\label{def:g10:stats:mean}
The \emph{mean}\index{mean} of the series is
\[
\bar x = \frac{x_1 + x_2 + \dots + x_N}{N}.
\]
If the values $v_1, \dots, v_k$ appear with counts $n_1, \dots, n_k$, the
same number is computed as the \emph{weighted mean}
\[
\bar x = \frac{n_1 v_1 + n_2 v_2 + \dots + n_k v_k}{n_1 + n_2 + \dots + n_k}.
\]
\end{definition}

\begin{example}\label{ex:g10:stats:mean}
For the quiz of \cref{ex:g10:stats:frequencies}:
\[
\bar x = \frac{1 \times 0 + 2 \times 1 + 4 \times 2 + 6 \times 3
+ 5 \times 4 + 2 \times 5}{20}
= \frac{0 + 2 + 8 + 18 + 20 + 10}{20}
= \frac{58}{20} = 2.9 .
\]
\end{example}

\begin{definition}[Median]\label{def:g10:stats:median}
The \emph{median}\index{median} is a value that splits the \emph{sorted}
series in two halves: at least half the values are $\leq$ the median, and
at least half are $\geq$ it. In practice, sort the $N$ values;
\begin{itemize}
  \item if $N$ is odd, the median is the middle value, in position
        $\frac{N+1}{2}$;
  \item if $N$ is even, take the \omterm{prop:g10:coordgeom:midpoint}{midpoint} of the two middle values, in
        positions $\frac N2$ and $\frac N2 + 1$.
\end{itemize}
\end{definition}

\begin{example}\label{ex:g10:stats:median}
Sorted series of $7$ house prices (in thousands): $120$, $150$, $160$,
$180$, $210$, $240$, $900$. The \omterm{def:g10:stats:median}{median} is the $4$th value, $180$. The \omterm{def:g10:stats:mean}{mean}
is $\frac{1960}{7} = 280$ --- larger than $6$ of the $7$ prices!
One extreme value ($900$) pulls the \omterm{def:g10:stats:mean}{mean} far more than the \omterm{def:g10:stats:median}{median}: the
\omterm{def:g10:stats:median}{median} is \emph{robust}, the \omterm{def:g10:stats:mean}{mean} is not.
\end{example}

\section{Measures of spread}

\begin{definition}[Range, quartiles]\label{def:g10:stats:quartiles}
For a sorted series:
\begin{itemize}
  \item the \emph{range}\index{range} is the difference between the
        largest and smallest values;
  \item the \emph{first quartile}\index{quartile} $Q_1$ is the smallest
        value such that at least a quarter ($25\%$) of the values are
        $\leq Q_1$; the \emph{third quartile} $Q_3$ is the smallest value
        such that at least three quarters ($75\%$) of the values are
        $\leq Q_3$;
  \item the \emph{interquartile range}\index{interquartile range} $Q_3 - Q_1$ measures the spread of
        the central half of the data.
\end{itemize}
\end{definition}

\begin{method}[Finding the quartiles]\label{met:g10:stats:quartiles}
For a series of $N$ sorted values:
\begin{enumerate}
  \item compute $\frac N4$; if it is not an \omterm{def:g10:numbers:sets}{integer} round it \emph{up} to
        the next \omterm{def:g10:numbers:sets}{integer}; $Q_1$ is the value in that position;
  \item compute $\frac{3N}{4}$; round up if needed; $Q_3$ is the value in
        that position.
\end{enumerate}
\end{method}

\begin{example}\label{ex:g10:stats:quartiles}
Take the $11$ sorted values
\[
2,\ 3,\ 5,\ 5,\ 6,\ 7,\ 8,\ 9,\ 9,\ 10,\ 12 .
\]
\omterm{def:g10:stats:median}{Median}: position $\frac{11+1}{2} = 6$, so the \omterm{def:g10:stats:median}{median} is $7$.
$Q_1$: $\frac{11}{4} = 2.75$, round up to $3$; the $3$rd value is $5$.
$Q_3$: $\frac{33}{4} = 8.25$, round up to $9$; the $9$th value is $9$.
\omterm{def:g10:stats:quartiles}{Range}: $12 - 2 = 10$; \omterm{def:g10:stats:quartiles}{interquartile range}: $9 - 5 = 4$.
\end{example}

\begin{omfigure}
\begin{tikzpicture}[scale=0.85]
  \draw[-Stealth] (1,0) -- (13.4,0) node[below, font=\small] {values};
  \foreach \x in {2,4,6,8,10,12}
    { \draw (\x,-0.07) -- (\x,0.07);
      \node[below, font=\footnotesize] at (\x,-0.1) {\x}; }
  \draw[omDef, thick, fill=omDef!12] (5,0.5) rectangle (9,1.3);
  \draw[omDef, very thick] (7,0.5) -- (7,1.3);
  \draw[omDef, thick] (2,0.9) -- (5,0.9);
  \draw[omDef, thick] (9,0.9) -- (12,0.9);
  \draw[omDef, thick] (2,0.7) -- (2,1.1);
  \draw[omDef, thick] (12,0.7) -- (12,1.1);
  \node[above, font=\small] at (5,1.32) {$Q_1$};
  \node[above, font=\small] at (7,1.32) {median};
  \node[above, font=\small] at (9,1.32) {$Q_3$};
  \node[above, font=\small] at (2,1.12) {min};
  \node[above, font=\small] at (12,1.12) {max};
\end{tikzpicture}

{\small The box plot of \cref{ex:g10:stats:quartiles}: the box spans the
interquartile \omterm{def:g10:numbers:interval}{interval} $\intcc{Q_1}{Q_3}$ and contains the central half of
the data; the whiskers reach the extreme values.}
\end{omfigure}

\begin{example}[Comparing two groups]\label{ex:g10:stats:comparison}
Two classes take the same test (scores out of $20$):
\begin{center}
\renewcommand{\arraystretch}{1.2}
\begin{tabular}{c|ccc}
 & \omterm{def:g10:stats:median}{median} & $Q_1$ & $Q_3$ \\
\hline
class A & $12$ & $10$ & $14$ \\
class B & $12$ & $6$ & $17$ \\
\end{tabular}
\end{center}
Same \omterm{def:g10:stats:median}{median}, very different spreads: class A is homogeneous (half its
scores in $\intcc{10}{14}$), class B is heterogeneous (its central half
spans $\intcc{6}{17}$). Indicators of center alone never tell the whole
story.
\end{example}

\begin{remark}
Neither the \omterm{def:g10:stats:mean}{mean} nor the \omterm{def:g10:stats:median}{median} is ``better'': the \omterm{def:g10:stats:mean}{mean} uses every value
(and is needed to compute totals), the \omterm{def:g10:stats:median}{median} resists extreme values. A
serious summary of a series gives at least one indicator of center
\emph{and} one of spread.
\end{remark}

\section{Exercises}

\begin{exercise}[$\star$]\label{exo:g10:stats:1}
Here are the numbers of books read in a year by $15$ students:
\[
0,\ 1,\ 1,\ 2,\ 2,\ 2,\ 3,\ 3,\ 4,\ 4,\ 5,\ 6,\ 7,\ 8,\ 12 .
\]
Compute the \omterm{def:g10:stats:mean}{mean}, the \omterm{def:g10:stats:median}{median} and the \omterm{def:g10:stats:quartiles}{range} of this series.
\end{exercise}

\begin{exercise}[$\star$]\label{exo:g10:stats:2}
A die was rolled $50$ times:
\begin{center}
\renewcommand{\arraystretch}{1.2}
\begin{tabular}{c|cccccc}
outcome & $1$ & $2$ & $3$ & $4$ & $5$ & $6$ \\
\hline
count & $7$ & $9$ & $8$ & $10$ & $6$ & $10$ \\
\end{tabular}
\end{center}
Compute the \omterm{def:g10:stats:series}{frequency} of each outcome and the \omterm{def:g10:stats:mean}{mean} outcome.
\end{exercise}

\begin{exercise}[$\star$]\label{exo:g10:stats:3}
Find the \omterm{def:g10:stats:median}{median}, $Q_1$ and $Q_3$ of the sorted series
\[
1,\ 2,\ 4,\ 4,\ 5,\ 6,\ 6,\ 7,\ 8,\ 9,\ 9,\ 11 \qquad (N = 12),
\]
and draw the corresponding box plot.
\end{exercise}

\begin{exercise}[$\star$]\label{exo:g10:stats:4}
The \omterm{def:g10:stats:mean}{mean} of $24$ test scores is $11.5$. A $25$th student takes the test
and scores $19$. What is the new \omterm{def:g10:stats:mean}{mean} of the class?
\end{exercise}

\begin{exercise}[$\star\star$]\label{exo:g10:stats:5}
A student has a \omterm{def:g10:stats:mean}{mean} of $12$ after $5$ tests, all with the same weight.
What score on the $6$th test would raise the \omterm{def:g10:stats:mean}{mean} to $13$? Is a \omterm{def:g10:stats:mean}{mean} of
$14$ reachable (scores are out of $20$)?
\end{exercise}

\begin{exercise}[$\star\star$]\label{exo:g10:stats:6}
Invent a series of $7$ nonnegative values whose \omterm{def:g10:stats:median}{median} is $10$ and whose
\omterm{def:g10:stats:mean}{mean} is larger than $20$; then one whose \omterm{def:g10:stats:median}{median} is $10$ and whose \omterm{def:g10:stats:mean}{mean} is
smaller than $6$. Why can the \omterm{def:g10:stats:mean}{mean} of such a series never be smaller than
$\frac{40}{7}$? What do these examples show?
\end{exercise}

\begin{exercise}[$\star\star$]\label{exo:g10:stats:7}
In a company, the $9$ employees earn $1800$ each per month and the
director earns $12000$.
\begin{enumerate}
  \item Compute the \omterm{def:g10:stats:mean}{mean} and the \omterm{def:g10:stats:median}{median} salary.
  \item Which indicator best describes a ``typical'' salary here? Why?
\end{enumerate}
\end{exercise}

\begin{exercise}[$\star\star$]\label{exo:g10:stats:8}
A class of $12$ boys has a \omterm{def:g10:stats:mean}{mean} height of $168$ cm; the $18$ girls of the
same class have a \omterm{def:g10:stats:mean}{mean} height of $163$ cm. Compute the \omterm{def:g10:stats:mean}{mean} height of the
whole class. (Careful: it is \emph{not} the \omterm{prop:g10:coordgeom:midpoint}{midpoint} of $168$ and $163$
--- weight the means by the group sizes.)
\end{exercise}

\begin{exercise}[$\star\star\star$]\label{exo:g10:stats:9}
A series of $N$ values $x_1, \dots, x_N$ has \omterm{def:g10:stats:mean}{mean} $\bar x$. Each value is
transformed into $y_i = a x_i + b$ for fixed numbers $a$ and $b$.
\begin{enumerate}
  \item Show that the \omterm{def:g10:stats:mean}{mean} of the new series is $a \bar x + b$.
  \item The temperatures of a week, in degrees Celsius, have \omterm{def:g10:stats:mean}{mean} $20$.
        What is their \omterm{def:g10:stats:mean}{mean} in degrees Fahrenheit
        ($F = 1.8\,C + 32$)?
\end{enumerate}
\end{exercise}
